% Modernised reading — fonds Alexandre Grothendieck, shelfmark 50, pages 1-10.
%
% This is an INTERPRETATION, not a transcription. Everything the critical
% apparatus carried moves into footnotes. In any dispute of substance the
% transcription governs: whoever cites Grothendieck must cite batch-01.fr.tex.
%
% Pass: Opus 5 (claude-opus-5), 2026-09-02 — first pass, unchecked against the
% pages by a human. The folder was read whole; it holds one batch, pages 1-10,
% of which page 1 is a wrapper and pages 3 and 5 an unrelated typescript.
%
% Notation fixed for the whole document. The manuscript's A* and D(A) are the
% same object — the dual — and it is written A^vee throughout, D being kept
% only for the duality functor of the axiomatic half. Nerm / N.S. is written NS.
% Pi is the group law and is written m. The two positivity cones are kept as he
% writes them, N^+ (the weak one) and N^> (the strong one), and the reading
% says once which is which. Xi is a class throughout, whether of a divisor
% (pages 2 and 4) or of the abstract N(A) (pages 6-10) — they are the same
% thing, and the whole point of the folder is that they are.
%
% Substantive departures from the pages, each footnoted where it occurs:
% — page 9 calls the Rosati map an "involution de l'anneau Hom(A,A)". It is an
%   ANTI-involution: (alpha beta)' = beta' alpha', as its own definition
%   phi^-1 t(-) phi shows. The reading says anti-involution.
% — page 8 says the sigma_i are forms of degree 2i in alpha. They are of degree
%   i: the page's own sigma_1 = tr is linear and its sigma_{2r} = nu is of
%   degree 2r. The reading says i and footnotes the page.
% — page 7's axiom (e) 2) states strict positivity over N^>(A) and then writes
%   its "i.e." with N^+(A). Only the strong cone can carry strict positivity —
%   a nef class of square zero exists on E x E — so the reading states it for
%   N^> and footnotes the page.
% — pages 6 to 10 lie in the folder in reverse. The reading reorders them
%   10, 9, 8 (from (d)), 7, then the head of 8, then 6, and footnotes the three
%   pieces of evidence: his own "apres d)" on page 7, the relative clause
%   opening page 9 which continues the boxed map closing page 10, and the fact
%   that (d) begins mid-sheet on page 8 with the space above it filled by what
%   continues the foot of page 7.
% — the antisymmetry of N(A) is his; the classical Neron-Severi group is the
%   group of SYMMETRIC homomorphisms A -> A^vee. The two differ by the sign of
%   the double-duality isomorphism kappa and both are legitimate; the reading
%   keeps his and footnotes the classical convention.
%
% What the folder announces and never establishes: the Corollaire of page 6,
% for which only the word is written — the positivity of the Rosati involution,
% which the two formulas above it give at once, and which the reading supplies
% as the edition's and not the page's; the ten-line diagonal note at the foot
% of page 8; the three words after "Le" on page 10; and the line of page 4 that
% carries "critere de Weil" and "variations", of which nothing else reads.
\documentclass[11pt,a4paper]{article}
% Le préambule commun à toutes les transcriptions du fonds Grothendieck.
%
% Il tient en une idée : **l'incertitude se compose, elle ne se résout pas.**
% Une transcription qui lisse un mot douteux a détruit la seule chose pour
% laquelle on la faisait — un lecteur doit pouvoir savoir, à chaque endroit,
% ce qui était sur le feuillet et ce qui a été ajouté. Les six macros qui
% suivent sont donc l'appareil critique, et non de la décoration.
%
% Les mêmes commandes sont reconnues par `scripts/render.mjs`, qui produit la
% vue de lecture du site. Une macro ajoutée ici doit l'être aussi là-bas,
% faute de quoi le rendu échouera bruyamment — ce qui est voulu : un rendu
% silencieusement faux est pire qu'un rendu absent.

\ProvidesPackage{grothendieck}[2026/08/08 Transcriptions du fonds Grothendieck]

\usepackage{amsmath,amssymb,amsthm}
\usepackage{mathrsfs}
\usepackage{stmaryrd}
% Les diagrammes commutatifs sont partout dans ce fonds : c'est la langue
% même de la géométrie algébrique de Grothendieck, et une transcription qui
% les rendrait en prose ne transcrirait rien.
\usepackage{tikz-cd}
% Français par défaut — c'est la langue des feuillets — mais l'anglais est
% chargé aussi : la traduction et le résumé sont des documents anglais, et la
% typographie française (espaces avant les deux-points) y serait une faute.
% Ces documents basculent par \selectlanguage{english} en tête de corps.
\usepackage[english,french]{babel}
\usepackage{fontspec}
\usepackage{geometry}
\geometry{a4paper,margin=25mm,marginparwidth=28mm}
\usepackage{marginnote}
\usepackage{xcolor}
% ulem, not soul. `soul`'s \st and \ul are built on a character-by-character
% scan that XeLaTeX's font machinery breaks: the document fails with an
% undefined control sequence rather than degrading. ulem does the same two jobs
% and is Unicode-engine safe. `normalem` stops it hijacking \emph, which the
% transcriptions use for Grothendieck's own emphasis.
\usepackage[normalem]{ulem}
\usepackage{hyperref}
\hypersetup{colorlinks=true,linkcolor=black,urlcolor=black,pdfborder={0 0 0}}

% --- Métadonnées du lot -----------------------------------------------------
% Reprises telles quelles par le script de rendu, qui les affiche en tête de
% la vue de lecture. Elles fixent ce que le fichier prétend couvrir : sans
% elles, une transcription flotte sans point d'attache dans un fonds de seize
% mille feuillets.

\newcommand{\folder}[1]{\gdef\grothFolder{#1}}
\newcommand{\batch}[1]{\gdef\grothBatch{#1}}
\newcommand{\leaves}[2]{\gdef\grothFirst{#1}\gdef\grothLast{#2}}
\newcommand{\foldertitle}[1]{\gdef\grothFoldertitle{#1}}
\gdef\grothFolder{?}\gdef\grothBatch{?}\gdef\grothFirst{?}\gdef\grothLast{?}\gdef\grothFoldertitle{}

% --- Le résumé --------------------------------------------------------------
%
% Il ouvre la lecture modernisée, sous le titre « Résumé », et il est composé
% en petit. Ce n'est pas un ornement : c'est le seul passage du document qui
% ne soit pas des mathématiques, et un lecteur doit pouvoir le reconnaître
% avant de l'avoir lu — puis le sauter s'il connaît déjà le sujet, ou s'y
% arrêter s'il ne le connaît pas. Le filet qui le suit marque la reprise du
% corps.

\newenvironment{resume}
  {\small\setlength{\parskip}{0.45em}}
  {\par\medskip\hrule\medskip}

% --- Mots-clés ---------------------------------------------------------------
%
% En anglais, en clôture du résumé de la lecture modernisée : le vocabulaire
% moderne sous lequel chercher ce que les feuillets construisent. Le manifeste
% du site (`npm run manifest`) les extrait pour étiqueter la cote entière —
% cette ligne est la source unique des tags, il n'y a pas de fichier à côté.
\newcommand{\keywords}[1]{\par\smallskip\noindent
  {\footnotesize\textsc{Keywords} --- \textit{#1}}\par}

% --- L'appareil critique ----------------------------------------------------

% Le numéro de feuillet, en marge. C'est l'ancre qui permet de régler un
% désaccord sur une formule en regardant *un* feuillet plutôt qu'une chemise
% de six cents. Le site s'en sert aussi pour tourner les pages du fac-similé
% quand on fait défiler la transcription.
\newcommand{\leaf}[1]{%
  \marginnote{\footnotesize\color{gray!70}\textsf{#1}}%
  \label{leaf:#1}\ignorespaces}

% Illisible : on ne devine pas. Un mot inventé qui se lit comme les autres est
% le pire résultat possible.
\newcommand{\ill}{\textcolor{red!60!black}{[\,\ldots\,]}}

% Lecture douteuse : le mot est donné, et signalé comme incertain.
\newcommand{\uncertain}[1]{\uline{#1}}

% Ajout de l'éditeur — une lettre manquante, un mot sous-entendu.
\newcommand{\add}[1]{\textcolor{blue!50!black}{[#1]}}

% Biffé par Grothendieck lui-même. Ce qu'il a rayé fait partie du document :
% une rature dit souvent où la pensée a changé de direction.
\newcommand{\struck}[1]{\sout{#1}}

% Note du transcripteur, sur l'état matériel du feuillet.
\newcommand{\note}[1]{\par\smallskip{\small\color{gray!60!black}\textit{[#1]}}\par\smallskip}

% Note marginale de Grothendieck, distincte du corps du texte.
\newcommand{\marginal}[1]{\par\smallskip\noindent\hspace{1em}%
  {\small\color{gray!60!black}#1}\par\smallskip}

% --- Titre ------------------------------------------------------------------

\AtBeginDocument{%
  \begin{center}
    {\large\bfseries Fonds Alexandre Grothendieck --- cote n\textsuperscript{o}~\grothFolder}\\[2pt]
    {\itshape\grothFoldertitle}\\[4pt]
    {\small feuillets \grothFirst--\grothLast\ (lot \grothBatch)}\\[2pt]
    {\footnotesize\color{gray!60!black}Transcription établie d'après le fac-similé numérisé
     par l'Université de Montpellier.\\
     Les lectures incertaines sont soulignées, les illisibles notées [\,\ldots\,],
     les ajouts entre crochets.}
  \end{center}
  \vspace{1em}}

\endinput


\folder{50}
\pages{1}{10}
\dating{[à partir de 1965]}
\watermark{Édition de démonstration\\interprétation personnelle de l'œuvre}
\foldertitle{Sur bouquin Lang : notes manuscrites (s.d.) — lecture modernisée du dossier entier}

\begin{document}

\section*{Résumé}

\begin{resume}
Quelqu'un lit un livre, et en cours de lecture entreprend de le récrire
autrement : non pas mieux, mais plus haut, de manière à voir ce qui, dans la
théorie, tient à la géométrie et ce qui n'y tient pas.

Le livre est celui de Serge Lang, \emph{Abelian Varieties}, paru en 1959 ;
Grothendieck l'appelle « le bouquin Lang », et c'est le titre qu'il a mis au
crayon sur la chemise. Une \emph{variété abélienne} est un objet géométrique
dont les points s'additionnent — la courbe elliptique en est le premier
exemple, et en toute dimension elles sont l'analogue algébrique du tore. À
chacune est attachée une seconde variété abélienne de même dimension, sa
\emph{duale}, qui paramètre les manières de tordre les fonctions sur la
première sans rien changer d'essentiel.

Le problème qui organise tout est celui-ci. Sur une telle variété on considère
une hypersurface, ou plutôt sa classe $\xi$ ; on la translate par un point $s$
et on compare la translatée à celle de départ. La comparaison définit, quand
$s$ varie, une application $\varphi_{\xi}$ de la variété vers sa duale. Deux
questions se posent aussitôt : quelles classes donnent l'application nulle, et
quelle structure porte l'ensemble de toutes les classes ?

La réponse à la première tient en une équation. Les classes invisibles pour
$\varphi$ sont exactement celles qui sont \emph{additives} : celles dont la
valeur en une somme de deux points est le produit des valeurs. Les quatre
premières lignes du feuillet 2 énoncent cette équivalence sous quatre formes —
invariance par translation, additivité, appartenance au $\mathrm{Pic}^{\tau}$,
provenance d'une extension de groupes par $\mathbb{G}_m$ — et notent en regard,
d'un mot chacun, les arguments qui les joignent. C'est un condensé, non une
démonstration : le feuillet ne prouve rien, il range.

La réponse à la seconde est l'objet du reste du dossier, et c'est là que la
lecture devient un travail. Une classe, s'aperçoit-on, n'est pas un objet
géométrique : c'est une \emph{forme alternée}. Une fois cela vu, presque tout
le formalisme du livre se détache de la géométrie. Il suffit d'une catégorie
additive munie d'une dualité — un procédé qui échange les objets et les
flèches, et qui redonne l'objet de départ quand on l'applique deux fois — pour
que la notion de forme alternée ait un sens ; il suffit de cette notion pour
que se décompose le groupe des classes d'un produit ; il suffit d'une
involution attachée à une classe non dégénérée pour retrouver ce que la théorie
classique appelle l'involution de Rosati. Rien de tout cela ne demande de
points, ni de corps de base, ni même de nombres.

Ce qui, en revanche, ne se laisse pas rendre formel, ce sont deux choses, et le
dossier les isole avec netteté en les posant comme axiomes plutôt qu'en les
démontrant : la \emph{positivité} — il existe un cône de classes « positives »,
stable par somme et par multiples entiers — et la \emph{forme d'intersection},
qui à $r$ classes associe un entier, $r$ étant la dimension. De ces deux
données descendent le degré d'un endomorphisme, sa trace, son polynôme
caractéristique, et pour finir le théorème vers lequel tout le dossier va :
une formule qui exprime la trace comme un rapport de nombres d'intersection.
Le corollaire qui en découle immédiatement — la positivité de l'involution de
Rosati, c'est-à-dire que $\mathrm{tr}(\alpha\alpha')$ est strictement positif
pour tout $\alpha$ non nul — est annoncé au dernier feuillet par son seul mot,
« Corollaire », suivi d'une page blanche.

Un mot au crayon, dans la marge du feuillet 9, dit tout le projet : en face de
la propriété qui décompose les classes d'un produit, il écrit « purement
formel ». C'est un jugement sur ce qu'il vient d'écrire, et c'est la ligne de
partage qu'il cherchait en ouvrant le livre. Lire, pour lui, c'est
apparemment tracer cette ligne : séparer ce qu'une théorie doit à son objet de
ce qu'elle doit à sa seule syntaxe, et n'appeler contenu que le premier. La
forme axiomatique où il aboutit — une catégorie autoduale, un cône positif, une
forme d'intersection à valeurs entières — est du reste celle qu'on rencontrera
plus tard ailleurs, dans les conjectures standard sur les motifs et dans les
structures de Hodge polarisées ; le dossier ne le dit pas, et n'a aucune raison
de le dire.

\keywords{abelian variety, dual abelian variety, Néron-Severi group, primitive line bundle, theorem of the cube, Barsotti-Weil formula, numerical equivalence, Matsusaka's theorem, Hopf-Borel theorem, self-dual additive category, alternating form, polarisation, ample cone, Rosati involution, intersection form, degree of an isogeny, characteristic polynomial of an endomorphism, trace formula, positivity}
\end{resume}

\pagerange{1}{10}
\section*{Le dossier, ses deux moitiés, et l'ordre des feuillets}

Dix feuillets. Le premier est une chemise grise portant, au crayon dans
l'angle, quatre mots de sa main : \emph{Sur bouquin Lang}. Les feuillets 3 et 5
sont le recto et le verso d'un tapuscrit Bourbaki sur les valuations, étranger
aux notes ; ils ne comptent pas.\footnote{L'inventaire de Montpellier prévient
que « les pages au verso des documents sont souvent sans rapport avec les notes
mathématiques ». C'est le cas ici, et c'est aussi ce qui date le papier plutôt
que les notes.} Restent huit feuillets, qui sont deux choses distinctes.

\begin{itemize}
\item Les feuillets 2 et 4 sont une même feuille de brouillon, au crayon et à
  l'encre noire, écrite en colonnes et non en lignes. Elle porte les quatre
  conditions équivalentes qui caractérisent les classes additives, les noms des
  arguments qui les joignent, la chaîne des trois équivalences entre classes de
  diviseurs, et l'encadrement de la dimension du Picard.
\item Les feuillets 6 à 10 sont une suite continue à l'encre bleue, qui pose
  l'axiomatique, lettre par lettre, de (a) à (e), puis en tire le théorème de
  la trace.
\end{itemize}

\subsection*{L'ordre de la suite}

Les feuillets 6 à 10 ne se lisent pas dans l'ordre où ils sont classés : ils y
sont, pour l'essentiel, à l'envers. L'ordre de lecture est 10, 9, puis le
feuillet 8 à partir de son milieu, puis 7, puis le haut du feuillet 8, puis
enfin 6.\footnote{Trois indices, aucun décisif seul, concordants ensemble. Le
feuillet 7 porte dans l'angle supérieur gauche, la feuille tournée d'un quart
de tour, « après d) » de sa main : c'est son propre repère, et il place ce
feuillet après celui qui porte l'axiome (d), lequel est le feuillet 8. Le
feuillet 9 s'ouvre sur une relative — « qui d'ailleurs s'annule sur
$\mathrm{Im}\,N(A) + \mathrm{Im}\,N(B)$ » — qui prolonge sans rupture
l'application encadrée sur laquelle se ferme le feuillet 10. Enfin l'axiome (d)
commence au milieu du feuillet 8, et ce qui occupe le haut de ce feuillet
prolonge exactement les dernières lignes du feuillet 7 : il a donc écrit (d)
en laissant le haut libre, est passé au feuillet 7 pour (e), et, celui-ci
rempli, est revenu remplir le haut resté blanc. Cette reconstitution est celle
de l'édition ; la transcription, elle, consigne les faits sans trancher.}

\subsection*{Les notations, fixées ici pour tout ce qui suit}

$A$, $B$ désignent des variétés abéliennes — ou, dans la moitié axiomatique,
des objets d'une catégorie additive $\mathcal{A}$. On écrit $A^{\vee}$ pour la
duale ; c'est le $A^{*}$ du feuillet 2 et le $D(A)$ des feuillets 6 à 10, et
c'est le même objet.\footnote{$D$ est gardé, mais seulement comme \emph{nom du
foncteur} de dualité, là où la moitié axiomatique en a besoin comme d'une
donnée et non comme d'un résultat.} On écrit $\mathrm{NS}$ pour le groupe de
Néron--Severi ; le manuscrit l'abrège « Nérm » au feuillet 2 et « N.S. » au
feuillet 4.\footnote{« Nérm » est la lecture littérale des quatre lettres ; la
flèche vers $\underline{\mathrm{Hom}}_{S\text{-}gr}(A, A^{\vee})$ ne laisse pas
d'autre possibilité que Néron--Severi, et le dossier 49, de la même main et du
même sujet, écrit « Nér.-Sév. ».} La loi de groupe est notée $m : A \times A
\to A$ ; le feuillet 2 l'appelle $\pi$. Enfin $\xi$ est une classe, qu'il
s'agisse d'un diviseur (feuillets 2 et 4) ou d'un élément du groupe abstrait
$N(A)$ (feuillets 6 à 10) : le dossier a précisément pour objet de montrer que
c'est la même chose.

Deux notations demandent d'être distinguées, parce qu'elles ne diffèrent que
par un exposant. $N^{+}(A)$ est le cône \emph{faible} et $N^{>}(A)$ le cône
\emph{fort} ; en langage classique, $N^{>}$ est le cône des classes amples et
$N^{+}$ celui des classes numériquement positives. C'est $N^{>}$ qui porte la
positivité stricte de la forme d'intersection et c'est de lui que sortent les
polarisations.\footnote{Le manuscrit repasse l'exposant du premier $N$ au
feuillet 8 et il s'y lit tour à tour $\tau$, $+$ et $\neq$ ; les lignes
suivantes distinguent nettement les deux. La lecture les tient pour deux cônes
distincts, ce que la suite du dossier confirme à chaque emploi.}

\pagerange{2}{2}
\section*{Les classes additives, et les quatre conditions (feuillet 2)}

Le feuillet part de l'application qui donne son sens à tout le reste :
\[
  \mathrm{NS}_{A/S} \longrightarrow
  \underline{\mathrm{Hom}}_{S\text{-}gr}(A, A^{\vee}),
  \qquad \xi \longmapsto \varphi_{\xi},
\]
où $\varphi_{\xi}(s) = t_{s}^{*}\xi - \xi$, $t_{s}$ étant la translation par
$s$. Dire que cette application est injective, c'est décrire son noyau au
niveau de $\mathrm{Pic}$ — et c'est exactement ce que font les quatre
conditions du bas du feuillet.\footnote{Le manuscrit écrit la flèche partant de
$\mathrm{NS}$ et non de $\mathrm{Pic}$, ce qui présuppose l'injectivité qu'il
va justifier plus bas. L'image, que le feuillet ne mentionne pas, est le groupe
des homomorphismes \emph{symétriques} $A \to A^{\vee}$.}

\textbf{Les quatre conditions.} \emph{Soit $\xi$ une classe de fibré en
droites sur $A$. Les conditions suivantes sont équivalentes :}

\begin{enumerate}
\item $\xi$ \emph{est dans le noyau de} $\mathrm{Pic}_{A/S} \to
  \underline{\mathrm{Hom}}(A, A^{\vee})$, \emph{c'est-à-dire}
  $t_{s}^{*}\xi = \xi$ \emph{pour tout point $s$} ;
\item $m^{*}\xi - \mathrm{pr}_{1}^{*}\xi - \mathrm{pr}_{2}^{*}\xi = 0$
  \emph{sur} $A \times A$ ;
\item $\xi \in \mathrm{Pic}^{\tau}(A/S)$ ;
\item $\xi$ \emph{est dans l'image de} $\mathrm{Ext}^{1}_{S}(A,
  \mathbb{G}_{m}) \to \mathrm{Pic}(A)$.
\end{enumerate}

Le manuscrit les numérote (i), (i$'$), (ii), (iii) — la seconde est une
variante de la première, et le prime le dit.\footnote{Les deux projections sont
écrites sans étoile sur la feuille, à la différence de $m^{*}$ ; la lecture
rétablit les étoiles, qui sont exigées par l'énoncé. L'indice de
$\mathrm{Ext}^{1}$ est tracé comme son $\xi$ et se lit $S$, la base, que
demande le sens.}

L'équivalence de (1) et (2) est ce que la feuille appelle « banal », et elle
l'est : le membre de gauche de (2) est le fibré de Mumford
$\Lambda(\xi) = m^{*}\xi - \mathrm{pr}_{1}^{*}\xi - \mathrm{pr}_{2}^{*}\xi$,
dont la restriction à $\{s\} \times A$ est précisément $t_{s}^{*}\xi - \xi$.
Une classe additive est une classe invariante par translation, et
réciproquement.

L'équivalence de (2) et (4) est la \emph{formule de Barsotti--Weil} :
$\mathrm{Ext}^{1}(A, \mathbb{G}_{m}) \simeq A^{\vee}$. L'identité (2) est le
cocycle : un fibré additif est un fibré en droites muni d'une loi de groupe
au-dessus de celle de $A$, c'est-à-dire une extension de $A$ par
$\mathbb{G}_{m}$.\footnote{La feuille l'attribue au « raisonnement
giraudique (de Serre » — le mot, que la transcription donne comme lecture
douteuse, est écrit en interligne, et la parenthèse n'est jamais refermée. L'argument de Serre est celui des
\emph{Groupes algébriques et corps de classes} ; l'adjectif renvoie
vraisemblablement à Jean Giraud et au langage des torseurs, qui est en effet
celui où l'équivalence est la plus courte. La transcription donne le mot comme
lecture douteuse, et le lien avec Giraud est de l'édition.}

L'équivalence de (1) et (3) est la seule qui demande un fait sur les variétés
abéliennes et non sur les fibrés : $\mathrm{Pic}^{\tau} = \mathrm{Pic}^{0}$,
c'est-à-dire que le groupe de Néron--Severi d'une variété abélienne est sans
torsion. Le feuillet 4 en donne la raison en trois symboles, et on la reprend
là.\footnote{La feuille écrit « (ii) $\Rightarrow$ (i) par la propriété
additive de $\mathrm{Pic}^{\tau}$ » et « (i) $\Rightarrow$ (ii) c'est le th. de
Barsotti--Cartier--Serre ». La première attribution est celle qu'on donnerait
aujourd'hui — c'est l'absence de torsion, obtenue au feuillet 4 ; la seconde
nomme, sous les trois auteurs qui l'ont établie indépendamment, le théorème qui
concerne en réalité la condition (4), et l'inclusion
$\mathrm{Pic}^{0} \subset \mathrm{Pic}^{\tau}$ qu'elle prétend justifier est
immédiate. La lecture donne les équivalences sous leur forme reçue et laisse à
la transcription les quatre lignes telles qu'elles sont.}

En marge de ces quatre lignes il note $t_{s}D \sim D$ : la même condition en
langage de diviseurs, celle sous laquelle on la rencontre dans le livre de
Lang.

\pagerange{4}{4}
\section*{Néron--Severi : trois équivalences, et l'encadrement de la dimension (feuillet 4)}

\subsection*{Pourquoi la torsion disparaît}

En haut de la colonne, trois symboles portent tout le poids : $dA = A$. La
multiplication par un entier $d \neq 0$ est surjective sur une variété
abélienne — c'est une isogénie. Il suit que $\mathrm{Hom}(A, A^{\vee})$ est
sans torsion : si $d\lambda = 0$, alors $\lambda$ s'annule sur $dA = A$. Comme
$\mathrm{NS}(A)$ s'injecte dans ce groupe d'homomorphismes, il est sans
torsion lui aussi, et $\mathrm{Pic}^{\tau}$ se réduit à
$\mathrm{Pic}^{0}$.\footnote{La ligne du manuscrit porte « additivité de
$\mathrm{Pic}^{\tau}$, on », puis un mot que la feuille ne livre pas, puis
« spécial : $dA = A$ ». La lecture supplée l'argument, qui est le seul que ces trois symboles
puissent porter et qui est celui dont le feuillet 2 a besoin.}

\subsection*{La chaîne des trois équivalences}

Le feuillet dresse ensuite, deux fois — une esquisse dans la colonne de gauche,
une version encadrée plus bas —, la chaîne des relations d'équivalence entre
classes de diviseurs :

\begin{tikzcd}
\text{équivalence algébrique} \arrow[d, Rightarrow] \\
\text{équivalence } \tau \arrow[d, Rightarrow] \\
\text{équivalence numérique}
\end{tikzcd}

Sur une variété projective quelconque les deux implications sont strictes en
général, et la seconde est un théorème : l'équivalence numérique et
l'équivalence $\tau$ coïncident.\footnote{C'est le théorème de Matsusaka
(1957). Le feuillet porte, entourée, la mention « dû à M. » suivie d'un nom
que la transcription ne tient pas pour assuré — neuf ou dix lettres finissant
en « -aka ». Le contenu du schéma désigne ce théorème-là quel que
soit le sort de la lecture.} Sur une variété abélienne, la torsion ayant
disparu, les trois relations n'en font qu'une : c'est ce que le feuillet vient
d'établir, et c'est pourquoi la chaîne est dessinée ici et pas ailleurs.

Une flèche montant vers « N.S. » porte trois lignes de commentaire dont il ne
reste que « critère de Weil », « variations » et « pour les VA (bouquin
Lang) ».\footnote{Rien d'autre ne se lit sur ces trois lignes ; la lecture ne
propose pas de reconstitution.}

\subsection*{L'encadrement de la dimension}

Enfin, entouré d'un trait, l'encadrement qui donne la dimension de la duale :
\[
  \dim A \;\leq\; \dim \mathrm{Pic}\,A \;\leq\;
  \dim H^{1}(A, \mathcal{O}_{A}) \;\leq\; \dim A .
\]
Les trois inégalités portent chacune une justification en accolade, et la
chaîne se referme sur elle-même : les trois dimensions sont égales.

\begin{itemize}
\item \emph{Théorie des diviseurs non dégénérés} : un diviseur non dégénéré $D$
  donne un $\varphi_{D} : A \to A^{\vee}$ de noyau fini, donc une isogénie sur
  son image ; la duale est au moins aussi grosse que $A$.
\item \emph{Définition de Picard} : l'espace tangent à l'origine du foncteur de
  Picard est $H^{1}(A, \mathcal{O}_{A})$, d'où la deuxième inégalité.
\item \emph{Hopf--Borel} : $H^{*}(A, \mathcal{O}_{A})$ est une algèbre de Hopf
  graduée anticommutative, donc l'algèbre extérieure de son premier étage ;
  comme elle est nulle au-delà du degré $\dim A$, on a
  $\dim H^{1} \leq \dim A$.
\end{itemize}

C'est mot pour mot l'argument du dossier 49, où le même encadrement sert à
démontrer que $\mathrm{Pic}^{0}$ d'un schéma abélien projectif est un schéma
abélien.\footnote{Dossier 49, feuillets 2 et 4 : « $\dim A^{\vee}_{0} \leq
\dim H^{1} \leq \dim A_{0}$ », avec les mêmes accolades. Les deux dossiers sont
manifestement contemporains, et celui-ci est la note de lecture dont l'autre
est la rédaction.}

Le feuillet se termine sur une ligne de trois symboles, $A^{**}$, $A$,
$A^{*}$, qui n'est pas développée : la biduale, l'objet, la duale.

\pagerange{10}{10}
\section*{L'axiomatique : une catégorie additive autoduale (feuillet 10)}

Ici commence la seconde moitié du dossier, et le changement d'altitude est
complet : il n'y a plus ni points, ni corps de base, ni diviseurs. En tête du
feuillet, au crayon : « livre de Lang / VA ».

\textbf{(a)} Une catégorie additive $\mathcal{A}$.

\textbf{(b)} Une \emph{autodualité} : un foncteur additif $D : \mathcal{A}^{\circ}
\to \mathcal{A}$ muni d'un isomorphisme $\kappa : \mathrm{id}_{\mathcal{A}}
\xrightarrow{\ \sim\ } D \circ D$ satisfaisant la condition de compatibilité
habituelle.\footnote{Cette condition est $D(\kappa_{A}) \circ \kappa_{D(A)} =
\mathrm{id}_{D(A)}$ ; le manuscrit dit « la condition de compatibilité
habituelle » et ne l'écrit pas. Elle est ce qui fait de la symétrie ci-dessous
une involution, et c'est à ce titre que la lecture la supplée.}

On pose alors
\[
  B(A,B) = \mathrm{Hom}(A, D(B)) \xrightarrow{\ s\ } \mathrm{Hom}(B, D(A))
  = B(B,A),
  \qquad s(u) = D(u) \circ \kappa_{B},
\]
qui est un bifoncteur contravariant en $A$ et en $B$, et où $s$ est un
isomorphisme involutif. Sur $B(A,A)$, $s$ est donc une involution, et il y a un
sens à parler d'éléments symétriques et antisymétriques. On écrira
${}^{t}u$ pour $s(u)$.

\textbf{(c)} Pour tout objet $A$, un sous-groupe
\[
  N(A) \subset B(A,A)
\]
formé d'éléments \emph{antisymétriques}, dépendant de $A$ de façon
contravariante — et il note en regard : \emph{$N$ n'est pas additif en
$A$}.\footnote{Le groupe de Néron--Severi classique est au contraire le groupe
des homomorphismes \emph{symétriques} $\lambda : A \to A^{\vee}$, c'est-à-dire
ceux qui vérifient $\lambda^{\vee} \circ \kappa_{A} = \lambda$. Les deux
conventions ne diffèrent que par le signe de $\kappa$ : remplacer $\kappa$ par
$-\kappa$ échange symétrie et antisymétrie sans toucher à la condition de
compatibilité, et les deux choix sont donc également légitimes. Le sien est
celui qui rend visible ce dont il s'agit — une forme alternée, celle que la
première classe de Chern définit sur $H_{1}$ — et c'est aussi celui qui impose
la forme matricielle ci-dessous.}

Pour $\xi \in N(A)$, on note $\varphi_{\xi} : A \to D(A)$ le morphisme
correspondant.

\subsection*{La forme d'une classe sur un produit}

Soit $\xi \in N(A \times B)$. Il définit $\varphi_{\xi} : A \times B \to D(A)
\times D(B)$, qui est antisymétrique — donc, dit-il, \emph{nécessairement} de
la forme
\[
  \begin{pmatrix}
    \alpha_{\xi} & {}^{t}\lambda_{\xi} \\
    -\lambda_{\xi} & \beta_{\xi}
  \end{pmatrix},
  \qquad \alpha_{\xi} \in N(A), \quad \beta_{\xi} \in N(B), \quad
  \lambda_{\xi} \in \mathrm{Hom}(A, D(B)) = B(A,B).
\]
Le mot « nécessairement » porte tout : l'antisymétrie de la matrice force les
deux blocs diagonaux à être antisymétriques et le bloc inférieur gauche à être
l'opposé du transposé du bloc supérieur droit. Le seul paramètre libre est
$\lambda_{\xi}$, et l'on obtient ainsi
\[
  \xi \longmapsto \lambda_{\xi} : \; N(A \times B) \longrightarrow
  \mathrm{Hom}(A, D(B)) = B(A,B),
\]
que le manuscrit encadre.

\pagerange{9}{9}
\section*{Ce qui est purement formel (feuillet 9)}

Cette application s'annule sur $\mathrm{Im}\,N(A) + \mathrm{Im}\,N(B)$ — les
classes venues d'un seul facteur n'ont pas de bloc hors-diagonale — et le
feuillet énonce aussitôt :

\textbf{Propriété fondamentale 1.} \emph{L'application $\xi \mapsto
\lambda_{\xi}$ induit un isomorphisme}
\[
  N(A \times B) \big/ \bigl(N(A) \times N(B)\bigr)
  \;\xrightarrow{\ \sim\ }\; \mathrm{Hom}(A, D(B)).
\]

En marge, au crayon et d'une autre pointe : « purement \emph{formel} ». Le
jugement est exact — la décomposition matricielle donne l'énoncé sans autre
ingrédient — et c'est le mot le plus important du dossier : il marque où
s'arrête ce que la syntaxe fournit. Il éclaire aussi, rétrospectivement, le
N.B. du feuillet 10 : si $N$ n'est pas additif en $A$, c'est de cette
non-additivité-là qu'il s'agit, et $\mathrm{Hom}(A, D(B))$ en est la mesure
exacte.\footnote{Traduit dans la théorie classique, c'est la décomposition
$\mathrm{NS}(A \times B) \simeq \mathrm{NS}(A) \oplus \mathrm{NS}(B) \oplus
\mathrm{Hom}(A, B^{\vee})$, qu'on démontre d'ordinaire par le théorème du
carré et le principe de la bascule. Ce que le feuillet fait voir, c'est qu'elle
ne doit rien à la géométrie.}

\textbf{Propriété fondamentale 2.} \emph{Pour tout $A$ il existe des
$\xi \in N(A)$ telles que $\varphi_{\xi} : A \to D(A)$ soit un
$\mathbb{Q}$-isomorphisme.}

C'est l'existence d'une polarisation, et ce n'est pas
formel : c'est ici que la géométrie rentre, sous forme d'axiome.\footnote{Trois
traits diagonaux barrent, à cet endroit, un « N.B. » entre crochets qui
invoquait les théorèmes des variétés abéliennes « pour $\mathbb{Q}$ », et dont
rien d'autre ne se lit. Sur une variété abélienne l'énoncé est le fait qu'il existe un fibré
ample, et $\varphi_{\xi}$ est alors une isogénie.}

\subsection*{L'involution de Rosati}

Une telle $\xi$ étant fixée, on pose
\[
  s_{\xi}(\alpha) = \varphi_{\xi}^{-1} \, {}^{t}\alpha \, \varphi_{\xi},
  \qquad \alpha \in \mathrm{Hom}(A,A)_{\mathbb{Q}},
\]
et l'on écrit $\alpha'$ pour $s_{\xi}(\alpha)$. C'est une
\emph{anti}-involution de l'anneau $\mathrm{Hom}(A,A)_{\mathbb{Q}}$ :
$(\alpha\beta)' = \beta'\alpha'$, comme la définition le montre
immédiatement.\footnote{Le manuscrit écrit « une involution de l'anneau
$\mathrm{Hom}(A,A)$ ». Le renversement des produits est une propriété de la
formule et non une nuance de vocabulaire : $\varphi^{-1}\,{}^{t}(\alpha\beta)\,
\varphi = \varphi^{-1}\,{}^{t}\beta\,{}^{t}\alpha\,\varphi =
(\varphi^{-1}\,{}^{t}\beta\,\varphi)(\varphi^{-1}\,{}^{t}\alpha\,\varphi)$.
C'est l'involution de Rosati.}

Il pose ensuite, pour $\alpha, \beta \in \mathrm{Hom}(A,A)_{\mathbb{Q}}$,
\[
  D_{\xi}(\alpha,\beta) = (\alpha+\beta)^{*}\xi - \alpha^{*}\xi -
  \beta^{*}\xi \;\in\; N(A),
\]
c'est-à-dire, en termes de $\varphi$,
\[
  \varphi_{D_{\xi}(\alpha,\beta)} =
  {}^{t}(\alpha+\beta)\,\varphi_{\xi}\,(\alpha+\beta)
  - {}^{t}\alpha\,\varphi_{\xi}\,\alpha
  - {}^{t}\beta\,\varphi_{\xi}\,\beta
  = {}^{t}\beta\,\varphi_{\xi}\,\alpha + {}^{t}\alpha\,\varphi_{\xi}\,\beta
  = \varphi_{\xi}\bigl(\alpha'\beta + \beta'\alpha\bigr).
\]
$D_{\xi}$ est donc, transportée par $\varphi_{\xi}$, la symétrisation du
produit pour l'involution de Rosati.\footnote{Dans $(\alpha+\beta)$ le signe
est tracé sur la feuille comme un petit anneau et non comme une croix ; c'est
l'identité finale qui impose la somme, et la transcription le note.} Trois
conséquences suivent, qui servent au feuillet 6 :
\[
  D_{\xi}(\alpha) := D_{\xi}(\alpha, \mathrm{id}_{A}) =
  \varphi_{\xi}(\alpha' + \alpha),
\]
\[
  D_{\xi}(\alpha,\beta) = D_{\xi}(\alpha'\beta) = D_{\xi}(\beta'\alpha),
  \qquad D_{\xi}(\alpha) = D_{\xi}(\alpha').
\]

\pagerange{8}{8}
\section*{La positivité, et les deux cônes (feuillet 8)}

\textbf{(d)} Pour tout objet $A$, on se donne dans $N(A)$ deux parties
\[
  N^{+}(A), \quad N^{>}(A) \;\subset\; N(A),
\]
dépendant fonctoriellement de $A$, et vérifiant :

\begin{enumerate}
\item $N^{+}(A) + N^{+}(A) \subset N^{+}(A)$ et $n\,N^{+}(A) \subset N^{+}(A)$
  pour $n > 0$ ; « kif kif » pour $N^{>}(A)$ — ce sont deux cônes ;
\item $N^{>}(A)$ engendre $N(A)$ ;
\item $\xi \in N^{>}(A) \Longrightarrow \varphi_{\xi}$ est un
  $\mathbb{Q}$-isomorphisme.
\end{enumerate}

La troisième clause est écrite deux fois : une première version, biffée, la
formulait à l'aide de la forme d'intersection — « $\xi \in N^{+}(A) \Rightarrow
I(\xi \ldots) > 0$ », puis « $\xi_{1}, \ldots, \xi_{r} \in N^{>}(A) \Rightarrow
I(\xi_{1}, \ldots, \xi_{r}) > 0$ » — avant qu'il ne la remplace par la
non-dégénérescence de $\varphi_{\xi}$ et ne renvoie la positivité de $I$ à
l'axiome (e).\footnote{La biffure est donc un déplacement d'axiome et non une
correction : la positivité passe de (d) à (e), où elle porte sur la forme
d'intersection et non plus sur les cônes seuls. C'est la seule trace, dans le
dossier, d'une hésitation sur l'architecture.} En langage classique : $N^{>}$
est le cône ample, et (3) dit qu'une classe ample définit une isogénie
$A \to A^{\vee}$, c'est-à-dire une polarisation.

Une note d'une dizaine de lignes court en diagonale au bas du feuillet, entre
deux longs traits. On y reconnaît $N^{+}(A)$, $n > 0$, le mot « diviseurs » et
$N^{>}(B)$ ; le reste ne se lit pas à la définition que porte le
film.\footnote{Elle est donc perdue pour cette édition. À en juger par ce qui
s'y reconnaît, elle discutait le contenu géométrique des deux cônes — ce que
« positif » veut dire pour un diviseur — ce qui est exactement la question que
les axiomes (d) et (e) contournent.}

\pagerange{7}{8}
\section*{La forme d'intersection, le degré, la trace (feuillets 7 et 8)}

\textbf{(e)} Pour tout objet $A$ on se donne une forme $r$-fois
\emph{linéaire}
\[
  (\xi_{1}, \ldots, \xi_{r}) \longmapsto I(\xi_{1}, \ldots, \xi_{r})
  \in \mathbb{Z},
\]
l'entier $r$ étant attaché à $A$ : $r = \dim(A)$.\footnote{Il écrit d'abord
$\mathbb{Q}$ puis le biffe pour $\mathbb{Z}$ ; la valeur d'un nombre
d'intersection est entière. La précision $r = \dim(A)$ est ajoutée en
interligne, reliée par un trait, et la transcription la donne comme lecture
douteuse.} C'est le nombre d'intersection de $r$ diviseurs sur une variété de
dimension $r$.

\textbf{2)} \emph{$I$ est strictement positive sur $N^{>}(A)$} : si
$\xi_{1}, \ldots, \xi_{r} \in N^{>}(A)$, alors $I(\xi_{1}, \ldots, \xi_{r}) >
0$.\footnote{Le manuscrit énonce la positivité « sur $N^{>}(A)$ » puis écrit
son « i.e. » avec $N^{+}(A)$. Seul le cône fort peut la porter : sur
$E \times E$ la classe de $\{0\} \times E$ est numériquement positive et de
carré nul. La lecture retient $N^{>}$ aux deux endroits.}

\textbf{3)} \emph{Pour $u : A \to B$ avec $\dim(A) = \dim(B)$, il existe un
entier $\nu(u)$ tel que}
\[
  I\bigl(u^{*}\xi_{1}, \ldots, u^{*}\xi_{r}\bigr)
  = \nu(u)\, I(\xi_{1}, \ldots, \xi_{r})
\]
\emph{pour toutes classes $\xi_{1}, \ldots, \xi_{r}$.} C'est le degré de $u$.
Il est bien déterminé dès qu'il existe une classe $\xi$ avec
$I(\xi, \ldots, \xi) \neq 0$, ce que la propriété 2) garantit.\footnote{La
phrase qui l'établit est en grande partie illisible : on y lit « ce qui le
détermine », « $\mathbb{Q}$-isomorphisme », « en particulier », et
« $\xi \in N(B)$ ». La lecture donne l'argument, qui est le seul disponible,
et le signale comme sien.}

En particulier, pour $\alpha$ un endomorphisme de $A$ :
\[
  \nu(\alpha) = \frac{1}{I(\xi, \ldots, \xi)}\,
  I\bigl(\alpha^{*}\xi, \ldots, \alpha^{*}\xi\bigr),
\]
qui est une forme de degré $2r$ en $\alpha$, à coefficients
entiers.\footnote{Le degré $2r$ se voit sur $\alpha = n \cdot \mathrm{id}$ :
$\nu(n) = n^{2r}$, puisque la multiplication par $n$ est de degré $n^{2r}$ en
dimension $r$. C'est aussi le rang de $H_{1}$.}

\subsection*{Le polynôme caractéristique}

De là, en développant :
\[
  \nu(\alpha + n \cdot 1_{A}) = n^{2r} + \sigma_{1}(\alpha)\,n^{2r-1} +
  \cdots + \sigma_{2r}(\alpha),
\]
avec $\sigma_{1}(\alpha) = \mathrm{tr}\,\alpha$ et $\sigma_{2r}(\alpha) =
\nu(\alpha)$, noté aussi $\det \alpha$. C'est, au signe des variables près, le
polynôme caractéristique de $\alpha$ agissant sur un espace de dimension $2r$,
et les $\sigma_{i}$ sont les fonctions symétriques élémentaires de ses valeurs
propres : $\sigma_{i}$ est donc une forme de degré $i$ en
$\alpha$.\footnote{Le feuillet 8 écrit « de degré $2i$ ». C'est un lapsus, et
le feuillet le contredit deux lignes plus haut : $\sigma_{1} = \mathrm{tr}$ est
linéaire et $\sigma_{2r} = \nu$ est de degré $2r$. L'homogénéité de
$\nu(\alpha + n \cdot 1)$, de degré $2r$ en $(\alpha, n)$, force le coefficient
de $n^{2r-i}$ à être de degré $i$.}

\subsection*{La trace comme rapport de nombres d'intersection}

Le haut du feuillet 8, qui prolonge ce qui précède, en tire la formule qui
justifie tout l'appareil :
\[
  \mathrm{tr}(\alpha) = \frac{r}{I(\xi, \ldots, \xi)}\,
  I\bigl(\xi, \ldots, \xi, D_{\xi}(\alpha)\bigr),
\]
où le second membre a $r-1$ arguments $\xi$ et un argument $D_{\xi}(\alpha) =
\varphi_{\xi}(\alpha' + \alpha)$, tandis que le dénominateur en a $r$. La trace,
définie plus haut comme un coefficient d'un polynôme, se calcule donc par un
rapport de nombres d'intersection.\footnote{Vérification sur $\alpha =
\mathrm{id}$, que le manuscrit ne fait pas : $D_{\xi}(\mathrm{id}) = 2\xi$,
donc le second membre vaut $2r$ — et $\nu(\mathrm{id} + n) = (n+1)^{2r}$ donne
bien $\sigma_{1}(\mathrm{id}) = 2r$. C'est la trace sur un espace de dimension
$2r$, non sur l'espace tangent.}

\pagerange{6}{6}
\section*{Le théorème de la trace, et le corollaire qui manque (feuillet 6)}

Le dernier feuillet énonce le théorème vers lequel les axiomes ont été taillés.

\textbf{Théorème.} \emph{Soit $\xi \in N^{>}(A)$ et soit $\alpha \mapsto
\alpha' = \varphi_{\xi}^{-1}\,{}^{t}\alpha\,\varphi_{\xi}$ l'involution de
Rosati de $\mathrm{Hom}(A,A)_{\mathbb{Q}}$. Alors}
\[
  \mathrm{tr}(\alpha\beta') = \frac{r}{I(\xi, \ldots, \xi)}\,
  I\bigl(\xi, \ldots, \xi, D_{\xi}(\alpha\beta')\bigr),
\]
\emph{d'où}
\[
  \mathrm{tr}(\alpha\alpha') = \frac{r}{I(\xi, \ldots, \xi)}\,
  I\bigl(\xi, \ldots, \xi, 2\,\alpha^{*}\xi\bigr).
\]

Le passage de la première formule à la seconde est l'identité (3) du feuillet 9
suivie d'un calcul d'une ligne : $D_{\xi}(\alpha,\alpha) = (2\alpha)^{*}\xi -
2\,\alpha^{*}\xi = 2\,\alpha^{*}\xi$, puisque $\xi$ est quadratique.

Puis, souligné, le mot \textbf{Corollaire} — et le feuillet s'arrête. Le reste
de la page est blanc, et le dossier avec elle.

Ce que les deux formules donnent immédiatement, et ce que le corollaire ne peut
guère être, est la \emph{positivité de l'involution de Rosati} : pour tout
$\alpha \neq 0$,
\[
  \mathrm{tr}(\alpha\alpha') > 0 .
\]
En effet $\alpha^{*}\xi$ appartient au cône positif dès que $\xi$ y appartient,
et l'axiome (e) 2) rend strictement positif le nombre d'intersection du second
membre.\footnote{Cette reconstitution est de l'édition : la feuille ne porte
que le mot. Elle est cependant contrainte de près — c'est l'énoncé pour lequel
la positivité a été posée en axiome au feuillet 8, c'est celui qui fait de
$\mathrm{Hom}(A,A)_{\mathbb{Q}}$ une algèbre semi-simple à involution positive,
et c'est celui qui, dans le livre de Lang, suit la formule de la trace. Le
lecteur qui voudrait la lettre du dossier ne trouvera, lui, que « Corollaire ».}

C'est là que le dossier s'interrompt, et l'endroit est parlant : tout ce qui
précède avait pour objet de rendre cet énoncé démontrable sans géométrie, et il
n'est pas écrit.

\end{document}
